\section{System Design}
\label{sec:design}

\subsection{Update Protocol}

Write requests are always funneled through the current coordinator. A non-coordinator replica
receiving a \texttt{WriteRequest} enqueues it locally (\texttt{m\_requested\_updates}, for later recovery) and forwards it to the coordinator; the coordinator enqueues every incoming write in
\texttt{m\_write\_queue} and processes it through \texttt{tryStartNextBroadcast()}, which starts a
new round only if no other update is currently in flight
(\texttt{m\_broadcast\_in\_progress}).

Each round assigns the write a timestamp $\langle e,i\rangle$ (current epoch, monotonically
increasing sequence number reset at every new epoch) and runs the classic two-phase protocol:

\begin{enumerate}[leftmargin=*]
  \item The coordinator broadcasts \texttt{UpdateMsg}$\langle e,i\rangle$ to all active replicas.
  \item Each replica stores the update as in-flight, replies with \texttt{AckMsg}, and arms a
  \texttt{WriteOKTimeout} to detect a coordinator crash mid-broadcast.
  \item Once the coordinator collects acknowledgements from a quorum
  $|Q| = \lfloor N/2 \rfloor + 1$ (itself included), it appends the update to its log and broadcasts
  \texttt{WriteOkMsg}.
  \item On \texttt{WriteOkMsg}, every non-coordinator replica appends the update to its own log
  and applies it to \texttt{PositionList}, replying to the client if it was the original request
  recipient. The coordinator does the same immediately after the broadcast loop, without going
  through \texttt{onWriteOkMsg} itself, since it never sends the message to itself.
\end{enumerate}

When the initiator of a write is not the coordinator, the coordinator additionally waits for an
\texttt{ApplyConfirmation} from it after sending \texttt{WriteOkMsg}, guarded by an
\texttt{ApplyTimeout}. This covers the case where the initiator crashes after the update was applied
by a quorum but before it manages to reply to the client: if the confirmation never arrives, the
coordinator sends the \texttt{WriteResult} to the client itself once the timeout fires.

\subsection{Reliable Broadcast Properties}

\begin{center}
\small
\begin{tabular}{|C{2.8cm}|C{13cm}|}
\hline
\textbf{Property} & \textbf{How it is satisfied} \\
\hline
Validity & A quorum of correct replicas is always available (system assumption), so the coordinator always eventually collects a quorum of ACKs and completes the
broadcast for any write issued by a correct client. \\
\hline
Integrity & Every update carries a unique $\langle e,i\rangle$ timestamp and is only ever applied
after being received via \texttt{WriteOkMsg} or \texttt{SynchronizationMsg}. The coordinator's single-broadcast-at-a-time implmentation together
with FIFO channels prevents duplicate delivery; during recovery, \texttt{UpdateLog.addLogIfAbsent}
explicitly guards against re-applying an update a replica already has. \\
\hline
Uniform Agreement & If the coordinator crashes after a quorum acknowledged an update but before every
replica received \texttt{WriteOkMsg}, the newly elected coordinator re-derives that update (from its
own in-flight state or its log) and re-disseminates it via \texttt{SynchronizationMsg} before serving
new writes, so all correct replicas converge to the same log. \\
\hline
Total Order & The coordinator keeps at most one update in flight at a time
(\texttt{m\_broadcast\_in\_progress}), assigning strictly increasing $\langle e,i\rangle$ timestamps
one round at a time; combined with FIFO channels, this guarantees every replica receives
\texttt{WriteOkMsg} messages in the same order they were generated. \\
\hline
\end{tabular}
\end{center}

\subsection{Sequential Consistency}

Channels are assumed reliable and FIFO, and \texttt{NetworkChannel} preserves this property between
any ordered pair of actors even with random delays. Clients are allowed to pipeline several requests
without waiting for a reply (\texttt{Client} tracks a list of outstanding reads and writes, not a
single one), so sequential consistency cannot rely on client-side blocking. It instead follows
directly from FIFO delivery combined with the fact that an actor processes its mailbox one message at
a time: requests from a given client to a given replica arrive, and are therefore handled, in the
exact order they were sent. Read requests are answered immediately from local state in that order,
while write requests are additionally serialized by the coordinator into a single totally-ordered
log, so any two clients reading the same index observe mutually consistent values.

\subsection{Crash Detection}

Three complementary timeouts detect coordinator failures at different protocol stages:

\begin{itemize}[leftmargin=*]
  \item \textbf{BroadcastTimeout}: armed by a replica after forwarding a write to the coordinator; it
  fires if the coordinator never starts the corresponding \texttt{UpdateMsg} broadcast.
  \item \textbf{WriteOKTimeout}: armed by a replica after acknowledging an \texttt{UpdateMsg}; it
  fires if the corresponding \texttt{WriteOkMsg} never arrives (coordinator crashed mid-broadcast).
  \item \textbf{HeartbeatReceiveTimeout}: armed by every non-coordinator replica and renewed on every
  received \texttt{HeartbeatRequest}; it fires when the coordinator is silently unreachable even with
  no pending write.
  \item 
  \textbf{ElectionGlobalTimeout}: armed when an election starts; it fires if the election does
  not complete in time (e.g.\ a ring participant or the winning candidate crashes mid-election),
  restarting the election from \texttt{m\_possible\_winner} if a winner was already identified, or
  from the old crashed coordinator otherwise.
\end{itemize}

Symmetrically, the coordinator uses per-replica \texttt{HeartbeatRequestTimeout}s to detect and evict
silently crashed replicas from its local \texttt{active\_replicas} map.

\subsection{Coordinator Election}

The ring is defined by sorting replica ids. When a coordinator crash is detected, election proceeds as follows:

\begin{enumerate}[leftmargin=*]
  \item The detecting replica enters \texttt{ELECTION} status and sends an \texttt{ElectionMsg}
  carrying a list of \texttt{CandidateEntry}$\langle$replicaId, lastKnownUpdate$\rangle$ to the next
  reachable replica in the ring.
  \item Each hop is acknowledged with \texttt{ElectionAckMsg} and guarded by
  \texttt{ElectionAckTimeout}: if a hop times out, the sender marks that replica as skipped
  (\texttt{m\_skipped\_in\_ring}) and forwards to the next one, so the election tolerates multiple
  consecutive crashes.
  \item A replica that sees its own id already in the candidate list has seen the full ring. It
  determines the winner as the candidate with the most recent $\langle e,i\rangle$, ties broken by
  highest id (\texttt{pickWinner}). Since a majority of replicas is always alive, at least one
  candidate is guaranteed to know the latest update.
  \item The winner promotes itself via \texttt{becomeCoordinator()}, starts a new epoch, and
  broadcasts \texttt{SynchronizationMsg} carrying any update that may not have reached every
  replica; other nodes adopt the new coordinator and epoch upon receiving it.
\end{enumerate}

Election completion is protected by \texttt{ElectionGlobalTimeout} (Section~\ref{sec:design},
Crash Detection). Separately, stale \texttt{ElectionMsg}s referring to an already-resolved election
are detected and dropped via the \texttt{m\_possible\_winner} bookkeeping, preventing the same
election from circulating indefinitely.

\subsection{Additional Assumptions}

\begin{itemize}[leftmargin=*]
  \item Heartbeats are implemented as request/response rather than a plain broadcast, so the
coordinator can also detect and evict silently crashed replicas, not just be detected by them.
\item Uniform Agreement guarantees an update is applied by every correct replica, but not that the
client is notified. The \texttt{ApplyConfirmation}/\texttt{ApplyTimeout} mechanism (Section~\ref{sec:write-confirmation}) covers this when
the initiator differs from the coordinator; the one gap left is the coordinator crashing between
reaching quorum and applying the update to itself while also being the initiator, where the client
may time out despite the write succeeding system-wide. Confirmed with the teaching assistants as
non-mandatory to handle since a full fix would require changes to classes shared with the provided tests.
\end{itemize}
